% §2.3 — Ammonia Adsorption/Desorption Process Dynamics
% Sources: Thesis/secs/3-governing_eqns/2-NH3ads_proc_dyn.tex
%          Thesis/secs/3-governing_eqns/2-1-sig_approx.tex
% ===========================================================
\subsection{Ammonia Adsorption/Desorption Process Dynamics \label{sec::sigma_deriv}}

% --- Discrete plug-flow figure ---
% \begin{figure}[!ht]
%     \centering
%     \includegraphics[width=\figWidth]{./figs/2-modeling/plug_flow_discrete.png}
%     \caption{Discrete plug-flow reactor model}
%     \label{fig::plug_flow_discrete}
% \end{figure}

The ammonia adsorption/desorption process involves all three reactions. Gaseous ammonia adsorbs onto free catalyst sites,
then either reacts with gaseous $\nox$ (Eley-Rideal mechanism) or decomposes via ammonia oxidation. Both processes free
adsorption sites for the next cycle. We first derive the dynamics when the catalyst remains unsaturated and then extend
the model to include saturation effects. The rate of change of adsorbed ammonia concentration is:
\begin{align}
    \dot{\con{\ammonia}}^{ads} &= \Gamma \underbrace{k_{ads} \con{\ammonia}^{in}}_{\gamma_{ads}}
        - \con{\ammonia}^{ads} \underbrace{\lr{k_{ads} \con{\ammonia}^{in} + k_{des} + k_{scr} \con{\nox}^{in}
        + k_{oxi}}}_{\gamma_{des}}
    \label{eqn::sigma_rate}
\end{align}

% --- Molar conservation at residence-time scale ---
\subsubsection{Molar Conservation at the Residence-Time Scale}
Let there be $n$ residence times within one sampling period:
\begin{align}
    t_s &= n \tau
    \label{eqn::ts_2_tau}
\end{align}
The molar storage on the catalyst surface changes at the end of every residence time as a fresh parcel of gaseous
reactants enters the chamber. Within the sample time $t_s$, the volumetric concentrations of the gaseous reactants are
assumed constant. The molar conservation for adsorbed ammonia can be written telescopically:
\begin{align*}
    \con{\ammonia}^{ads}(k + n\tau) &= \con{\ammonia}^{ads} (k)
        + \underbrace{\sum_{i=0}^{n-1} \int_{0}^{\tau} \dot{\con{\ammonia}}^{ads} (k + i\tau) \, dt}_{\Omega(k)}
\end{align*}
where $\Omega(k)$ is the total change in the surface concentration of adsorbed ammonia within one sample time.
Evaluating the integral under the assumption that the rate is constant within each residence time:
\begin{align}
    \Omega(k) &= n \tau \lr{\Gamma \gamma_{ads} (k) - \sigma(k) \gamma_{des} (k)}
    \label{eqn::omega}
\end{align}
where
\begin{align}
    \sigma(k) &= \frac{1}{n} \sum_{i=0}^{n-1} \con{\ammonia}^{ads} (k + i \tau)
    \label{eqn::sigma_def}
\end{align}
is the \textit{average surface concentration} of adsorbed ammonia over the sampling period. This quantity is unobservable.

% --- Zero-order hold approximation ---
\subsubsection{Zero-Order Hold Approximation}
The average surface concentration $\sigma(k)$ is approximated as the surface concentration at the beginning of the
sample time: $\sigma(k) \approx \con{\ammonia}^{ads}(k)$. As $n\tau = t_s$ is constant, the discrete update equation
becomes:
\begin{align}
    \sigma^{ub}(k+1) &= \sigma(k) + t_s k_{ads} \con{\ammonia}^{in} (k) \lr{\Gamma - \sigma(k)} \notag\\
        & \qquad - t_s \lr{k_{oxi} + k_{des}} \sigma(k) - t_s k_{scr} \con{\nox}^{in}(k) \sigma(k)
    \label{eqn::sigma_ub}
\end{align}
where the superscript $ub$ denotes the \textit{unbounded} (pre-saturation) update. The individual terms represent:
adsorption onto free sites, ammonia oxidation and desorption, and $\nox$ reduction, respectively.
